\documentclass[12pt,a4paper,french]{report}

\usepackage[utf8]{inputenc}
\usepackage[T1]{fontenc}
\usepackage{babel}
\usepackage{geometry}
\usepackage{xcolor}
\usepackage{hyperref}
\usepackage{graphicx}
\usepackage{listings}
\usepackage{fancyhdr}
\usepackage{titlesec}
\usepackage{tcolorbox}
\usepackage{enumitem}
\usepackage{booktabs}
\usepackage{longtable}
\usepackage{array}
\usepackage{tabularx}
\usepackage{float}
\usepackage{tikz}
\usepackage{amsmath}
\usepackage{setspace}
\usepackage{parskip}
\usepackage{charter}

% ─── Géométrie ───────────────────────────────────────────────
\geometry{
  top=2.5cm, bottom=2.5cm,
  left=2.5cm, right=2.5cm,
  headheight=15pt
}

% ─── Couleurs ────────────────────────────────────────────────
\definecolor{primaryblue}{RGB}{31,78,121}
\definecolor{lightblue}{RGB}{214,228,240}
\definecolor{accentblue}{RGB}{70,130,180}
\definecolor{darkgray}{RGB}{64,64,64}
\definecolor{lightgray}{RGB}{245,245,245}
\definecolor{successgreen}{RGB}{0,128,0}
\definecolor{warningorange}{RGB}{204,102,0}
\definecolor{errored}{RGB}{180,30,30}
\definecolor{codebg}{RGB}{248,248,248}

% ─── Hyperref ────────────────────────────────────────────────
\hypersetup{
  colorlinks=true,
  linkcolor=primaryblue,
  urlcolor=accentblue,
  citecolor=primaryblue,
  pdftitle={Documentation Pipeline Panel Admin},
  pdfauthor={Data Engineering Team},
}

% ─── En-têtes / pieds de page ────────────────────────────────
\pagestyle{fancy}
\fancyhf{}
\fancyhead[L]{\small\color{darkgray} Pipeline Panel Admin — Documentation Technique}
\fancyhead[R]{\small\color{darkgray} \leftmark}
\fancyfoot[C]{\small\color{darkgray} \thepage}
\renewcommand{\headrulewidth}{0.4pt}
\renewcommand{\footrulewidth}{0pt}

% ─── Titres de chapitres ─────────────────────────────────────
\titleformat{\chapter}[block]
  {\normalfont\huge\bfseries\color{primaryblue}}
  {\thechapter.}{1em}{}
  [\vspace{0.5ex}\titlerule[1pt]]

\titleformat{\section}[block]
  {\normalfont\Large\bfseries\color{primaryblue}}
  {\thesection}{1em}{}

\titleformat{\subsection}[block]
  {\normalfont\large\bfseries\color{accentblue}}
  {\thesubsection}{1em}{}

\titleformat{\subsubsection}[block]
  {\normalfont\normalsize\bfseries\color{darkgray}}
  {\thesubsubsection}{1em}{}

% ─── Boîtes colorées ─────────────────────────────────────────
\tcbuselibrary{skins,breakable}

\newtcolorbox{infobox}[1][]{
  enhanced, breakable,
  colback=lightblue!30, colframe=primaryblue,
  fonttitle=\bfseries\color{primaryblue},
  title={#1}, left=6pt, right=6pt, top=4pt, bottom=4pt
}

\newtcolorbox{warningbox}[1][Attention]{
  enhanced, breakable,
  colback=orange!10, colframe=warningorange,
  fonttitle=\bfseries\color{warningorange},
  title={#1}, left=6pt, right=6pt, top=4pt, bottom=4pt
}

\newtcolorbox{codebox}[1][]{
  enhanced, breakable,
  colback=codebg, colframe=darkgray!40,
  fonttitle=\bfseries\small\color{darkgray},
  title={#1}, left=6pt, right=6pt, top=4pt, bottom=4pt,
  fontupper=\small\ttfamily
}

% ─── Listings ────────────────────────────────────────────────
\lstset{
  backgroundcolor=\color{codebg},
  basicstyle=\small\ttfamily,
  breaklines=true,
  captionpos=b,
  commentstyle=\color{successgreen},
  keywordstyle=\color{primaryblue}\bfseries,
  stringstyle=\color{errored},
  numberstyle=\tiny\color{darkgray},
  numbers=left,
  stepnumber=1,
  numbersep=8pt,
  frame=single,
  framesep=4pt,
  rulecolor=\color{darkgray!40},
  showstringspaces=false,
  tabsize=4,
  language=Python,
  literate={é}{{\'e}}1 {è}{{\`e}}1 {ê}{{\^e}}1 {à}{{\`a}}1 {â}{{\^a}}1
           {î}{{\^i}}1 {ô}{{\^o}}1 {ù}{{\`u}}1 {û}{{\^u}}1
           {É}{{\'E}}1 {È}{{\`E}}1 {Â}{{\^A}}1 {Ô}{{\^O}}1
}

% ─── Commandes utilitaires ───────────────────────────────────
\newcommand{\scriptname}[1]{\texttt{\textbf{#1}}}
\newcommand{\bucket}[1]{\texttt{s3://#1}}
\newcommand{\colname}[1]{\texttt{#1}}
\newcommand{\etape}[2]{\textbf{Étape #1} — #2}

% ─── Espacement ──────────────────────────────────────────────
\setlength{\parskip}{6pt}
\setlength{\parindent}{0pt}
\onehalfspacing

% ─────────────────────────────────────────────────────────────
\begin{document}

% ─── Page de titre ───────────────────────────────────────────
\begin{titlepage}
  \centering
  \vspace*{2cm}
  {\huge\bfseries\color{primaryblue}
    Pipeline de Données\\[0.4em]
    Panel Administratif\\[0.4em]
    2015--2025
  \par}
  \vspace{1.5cm}
  \textcolor{accentblue}{\rule{10cm}{2pt}}
  \vspace{1cm}

  {\Large\bfseries Documentation Technique Complète\par}
  \vspace{0.5cm}
  {\large Version 2.0 — Architecture Python / Polars / MinIO\par}
  \vspace{2cm}

  \begin{tcolorbox}[colback=lightblue!40, colframe=primaryblue, width=12cm]
    \centering
    \textbf{Périmètre :} Indicateurs de la masse salariale des agents de l'État\\[4pt]
    \textbf{Couverture temporelle :} Janvier 2015 -- Décembre 2025\\[4pt]
    \textbf{Technologie :} Python · Polars · MinIO (S3) · openpyxl
  \end{tcolorbox}

  \vfill
  {\small\color{darkgray} Document généré le \today\par}
\end{titlepage}

% ─── Table des matières ──────────────────────────────────────
\tableofcontents
\newpage

% ═════════════════════════════════════════════════════════════
% CHAPITRE 1 — RÉSUMÉ EXÉCUTIF
% ═════════════════════════════════════════════════════════════
\chapter{Résumé exécutif}

Ce document décrit le pipeline de traitement de données qui produit les
indicateurs statistiques sur les salaires des agents de la fonction publique
à partir des fichiers de solde mensuels bruts (2015--2025).

\section{Vue d'ensemble}

Le pipeline transforme environ 120 fichiers Excel mensuels bruts en un panel
individu $\times$ période exploitable, puis calcule les indicateurs de
rémunération diffusés dans un fichier Excel structuré.

\begin{infobox}[Architecture en médaillon (Medallion Architecture)]
  \begin{itemize}[nosep]
    \item \textbf{Staging} : dépôt initial des fichiers sources (Excel mensuels +
          référentiels)
    \item \textbf{Bronze} : données brutes converties en Parquet, tout en
          \texttt{String}, sans transformation
    \item \textbf{Silver} : données enrichies, types corrects, situations
          normalisées, codes ANSTAT rattachés
    \item \textbf{Gold} : agrégats (masse salariale, effectifs) + panel
          individu $\times$ période avec grades imputés
  \end{itemize}
\end{infobox}

\section{Livrables produits}

\begin{enumerate}
  \item \textbf{Panel complet} \bucket{gold/panel\_admin/panel\_complet.parquet}
        et fichiers annuels \texttt{panel\_YYYY.parquet} : une ligne par
        agent $\times$ mois, avec grade imputé, code CITP, barème salarial.

  \item \textbf{Tables Gold d'agrégats} :
        \begin{itemize}[nosep]
          \item \bucket{gold/panel\_admin/masse\_salariale.parquet}
          \item \bucket{gold/panel\_admin/effectifs.parquet}
        \end{itemize}

  \item \textbf{Fichier d'indicateurs} : \texttt{indicateurs\_\{mode\}\_\{periode\}\_\{zeros\}.xlsx}
        (multi-feuilles : grades, CITP, barèmes, multi-postes, modèle ANSTAT)

  \item \textbf{Fichier avec sommaire navigable} : \texttt{indicateurs\_*\_SOMMAIRE.xlsx}
        (5 sections cliquables, feuilles de détail par grade/CITP/sexe/barème)

  \item \textbf{Rapports de qualité} dans \bucket{staging/panel\_admin/validation/} :
        taux de matching ANSTAT, doublons, colonnes suspectes.
\end{enumerate}

\section{Référentiels Staging}

Trois types de fichiers de référence sont déposés dans
\bucket{staging/panel\_admin/references/} avant toute exécution du pipeline :

\begin{enumerate}
  \item \textbf{FICHIER\_ANSTAT\_CODE.xlsx} — dictionnaire de codes administratifs
        (organisme, grade, emploi, affectation, service, fonction, poste) avec
        libellés normalisés pour le matching textuel.
  \item \textbf{CORRESPONDANCE\_GRILLE\_EMPLOI.xlsx} \textbf{(nouveau)} — table de
        correspondance emploi $\rightarrow$ barème / grille salariale, utilisée pour
        le fuzzy join de l'étape 07.
  \item Fichiers Excel mensuels déposés dans
        \bucket{staging/panel\_admin/fichiers\_mensuels/}.
\end{enumerate}

\section{Suppressions de lignes dans le pipeline}

Le pipeline supprime des lignes à deux étapes clés :

\begin{itemize}
  \item \textbf{Étape 03 (Bronze $\rightarrow$ Silver)} : toute ligne dont
        \colname{situation\_normalisee} vaut \texttt{autre} est supprimée.
        Les situations conservées sont \texttt{en\_activite},
        \texttt{regul\_indemnites}, \texttt{demi\_solde}. Les données F2
        associées à ces lignes sont également écartées.
  \item \textbf{Étape 07 (Calcul indicateurs)} : les lignes dont
        \colname{GRADE\_SOURCE} est \texttt{NF\_NON\_IMPUTABLE} sont exclues
        des calculs. Les agents de sexe inconnu sont aussi exclus.
\end{itemize}

% ═════════════════════════════════════════════════════════════
% CHAPITRE 2 — CONTEXTE ET OBJECTIFS
% ═════════════════════════════════════════════════════════════
\chapter{Contexte et objectifs}

\section{Contexte}

La production des indicateurs de la masse salariale de la fonction publique
repose sur des fichiers de solde mensuels (format Excel, deux feuilles :
données agents et codes de rémunération). Ces fichiers, produits depuis 2015,
présentent des hétérogénéités structurelles importantes : noms de colonnes
variables, lignes d'en-tête décalées, formats de date divers.

Le pipeline automatise l'intégralité de la chaîne de traitement, de
l'ingestion brute jusqu'à la production des indicateurs statistiques.

\section{Objectifs techniques}

\begin{itemize}
  \item Centraliser tous les fichiers sources sur un stockage objet partagé
        (MinIO) pour permettre le travail collaboratif depuis JupyterHub.
  \item Standardiser les données sur l'ensemble de la série temporelle malgré
        les hétérogénéités structurelles.
  \item Constituer un panel longitudinal individu $\times$ période de qualité,
        avec grades validés et codes CITP rattachés.
  \item Produire des indicateurs salariaux fiables (P25, P50, P75, min, max,
        moyenne) par grade, CITP, sexe, barème salarial.
  \item Générer un fichier Excel navigable destiné aux utilisateurs métier.
\end{itemize}

\section{Contraintes de conception}

\begin{itemize}
  \item \textbf{Mémoire limitée} : le panel complet (120+ fichiers, plusieurs
        millions de lignes) est traité en streaming via
        \texttt{scan\_parquet} / \texttt{sink\_parquet} de Polars.
  \item \textbf{Robustesse} : chaque étape gère les erreurs fichier par fichier.
        Un fichier en erreur n'arrête pas le traitement des autres.
  \item \textbf{Traçabilité} : chaque ligne du panel conserve son
        \colname{fichier\_source} et son \colname{GRADE\_SOURCE}
        (\texttt{OBSERVE}/\texttt{IMPUTE}/\texttt{NF\_NON\_IMPUTABLE}).
  \item \textbf{Reproductibilité} : les paramètres de calcul (période, mode
        salaire, seuil de winsorisation) sont centralisés en tête de chaque
        script.
\end{itemize}

% ═════════════════════════════════════════════════════════════
% CHAPITRE 3 — ARCHITECTURE TECHNIQUE
% ═════════════════════════════════════════════════════════════
\chapter{Architecture technique}

\section{Stack technologique}

\begin{tabularx}{\textwidth}{lX}
  \toprule
  \textbf{Composant} & \textbf{Rôle} \\
  \midrule
  Python 3.10+ & Langage de traitement \\
  Polars & Manipulation de DataFrames (plus rapide que pandas, API Lazy) \\
  boto3 & Client S3 pour MinIO \\
  fastexcel & Lecture Excel rapide (backend Polars) \\
  openpyxl & Écriture Excel stylisée (sortie finale) \\
  python-dotenv & Gestion des variables d'environnement \\
  MinIO & Stockage objet S3-compatible (auto-hébergé) \\
  JupyterHub & Environnement d'exécution partagé \\
  \bottomrule
\end{tabularx}

\section{Organisation du stockage objet}

\begin{codebox}[Structure des buckets MinIO]
staging/
  panel\_admin/
    fichiers\_mensuels/       ← Fichiers Excel sources (MMYYYY.xlsx)
    references/
      FICHIER\_ANSTAT\_CODE.xlsx
      CORRESPONDANCE\_GRILLE\_EMPLOI.xlsx   ← NOUVEAU
    pre\_analyse/             ← Rapports 01 (CSV)
    validation/              ← Rapports 04 + 03 (CSV qualité)
    exports\_gold/            ← Fichiers Excel finaux + CSV exports

bronze/
  panel\_admin/
    MMYYYY.parquet            ← Un fichier par mois (ex: 012015.parquet)

silver/
  panel\_admin/
    YYYY.parquet              ← Un fichier par année (ex: 2020.parquet)

gold/
  panel\_admin/
    masse\_salariale.parquet  ← Agrégats masse salariale
    effectifs.parquet         ← Agrégats effectifs
    panel\_YYYY.parquet       ← Panel annuel (ex: panel\_2020.parquet)
    panel\_complet.parquet    ← Panel toutes années fusionné
\end{codebox}

\section{Flux de données}

\begin{tikzpicture}[
  node distance=1.5cm,
  box/.style={rectangle, draw=primaryblue, fill=lightblue!30,
              minimum width=3cm, minimum height=0.8cm,
              text centered, font=\small\bfseries},
  arrow/.style={->, thick, color=primaryblue}
]
  \node[box] (staging) {STAGING\\(Excel)};
  \node[box, right=of staging] (bronze) {BRONZE\\(Parquet brut)};
  \node[box, right=of bronze] (silver) {SILVER\\(Parquet enrichi)};
  \node[box, right=of silver] (gold) {GOLD\\(Panel + agrégats)};
  \node[box, below=of gold] (xlsx) {Excel\\Indicateurs};

  \draw[arrow] (staging) -- (bronze) node[midway, above, font=\tiny] {02};
  \draw[arrow] (bronze)  -- (silver) node[midway, above, font=\tiny] {03};
  \draw[arrow] (silver)  -- (gold)   node[midway, above, font=\tiny] {05/06};
  \draw[arrow] (gold)    -- (xlsx)   node[midway, right, font=\tiny] {07/08};
\end{tikzpicture}

\section{Connexion à MinIO}

Chaque script supporte deux modes de connexion :

\begin{codebox}[Configuration locale vs JupyterHub]
# Depuis la machine locale
MINIO\_ENDPOINT   = "http://192.168.1.230:30137"
MINIO\_ACCESS\_KEY = "datalab-team"
MINIO\_SECRET\_KEY = "minio-datalabteam123"

# Depuis JupyterHub (décommenter)
# MINIO\_ENDPOINT = "http://minio.mon-namespace.svc.cluster.local:80"
\end{codebox}

Les scripts des étapes 03 à 07 utilisent des paramètres de retry renforcés :
5 tentatives maximum, mode adaptatif, \texttt{connect\_timeout=30s},
\texttt{read\_timeout=120s}.

% ═════════════════════════════════════════════════════════════
% CHAPITRE 4 — DESCRIPTION DES SCRIPTS
% ═════════════════════════════════════════════════════════════
\chapter{Description des scripts}

% ─────────────────────────────────────────────────────────────
\section{00\_upload\_staging.py — Dépôt initial des sources}
% ─────────────────────────────────────────────────────────────

\begin{warningbox}[À exécuter une seule fois]
  Ce script doit être lancé \textbf{une seule fois} avant toute exécution du
  pipeline. Il transfère les fichiers sources depuis le poste local vers
  MinIO staging, et crée les quatre buckets s'ils n'existent pas.
\end{warningbox}

\subsection{Objectif}

Déposer les fichiers sources sur MinIO afin que toute l'équipe puisse
travailler depuis staging, sans dépendance aux chemins locaux.

\subsection{Variables configurables}

\begin{tabularx}{\textwidth}{lX}
  \toprule
  \textbf{Variable} & \textbf{Description} \\
  \midrule
  \texttt{DOSSIER\_EXCEL} &
    Chemin local vers le dossier contenant les fichiers Excel mensuels
    (parcours récursif \texttt{**/*.xlsx}).
    Défaut : \texttt{consolidation\_solde\_2015\_2025/01\_data\_sources/fichiers\_solde\_mensuels} \\
  \texttt{FICHIER\_ANSTAT} &
    Chemin local vers le fichier ANSTAT\_CODE.
    Défaut : \texttt{.../FICHIER\_ANSTAT\_CODE\_2025.xlsx} \\
  \texttt{FICHIER\_CORRESPONDANCE} &
    \textbf{(Nouveau)} Chemin local vers la table de correspondance
    emploi $\rightarrow$ barème.
    Défaut : \texttt{.../CORRESPONDANCE\_GRILLE\_EMPLOI.xlsx} \\
  \texttt{MINIO\_ENDPOINT} & URL du serveur MinIO \\
  \texttt{MINIO\_ACCESS\_KEY} & Clé d'accès MinIO \\
  \texttt{MINIO\_SECRET\_KEY} & Clé secrète MinIO \\
  \bottomrule
\end{tabularx}

\subsection{Comportement détaillé}

\textbf{Création automatique des buckets.} Le script vérifie l'existence des
quatre buckets (\texttt{staging}, \texttt{bronze}, \texttt{silver},
\texttt{gold}) via \texttt{head\_bucket} et les crée si nécessaire.

\textbf{Upload des fichiers mensuels.} Tous les fichiers \texttt{.xlsx}
trouvés récursivement dans \texttt{DOSSIER\_EXCEL} sont uploadés vers
\bucket{staging/panel\_admin/fichiers\_mensuels/\{nom\_fichier\}}.
Un compteur OK/erreurs est affiché à la fin.

\textbf{Upload du fichier ANSTAT.} Le fichier \texttt{FICHIER\_ANSTAT} est
uploadé vers \bucket{staging/panel\_admin/references/\{nom\}}.

\textbf{Upload de la table de correspondance.} \textbf{(Nouveau)}
\texttt{CORRESPONDANCE\_GRILLE\_EMPLOI.xlsx} est uploadé vers
\bucket{staging/panel\_admin/references/CORRESPONDANCE\_GRILLE\_EMPLOI.xlsx}.
Si le fichier est absent localement, un message indique le chemin attendu
sans arrêter le script.

\subsection{Dépendances Python}

\begin{codebox}[]
pip install boto3 python-dotenv
\end{codebox}

\textit{Note : Polars n'est pas requis pour cette étape.}

\subsection{Fichiers produits}

\begin{itemize}[nosep]
  \item \bucket{staging/panel\_admin/fichiers\_mensuels/*.xlsx}
  \item \bucket{staging/panel\_admin/references/FICHIER\_ANSTAT\_CODE.xlsx}
  \item \bucket{staging/panel\_admin/references/CORRESPONDANCE\_GRILLE\_EMPLOI.xlsx}
\end{itemize}

% ─────────────────────────────────────────────────────────────
\section{01\_pre\_analyse.py — Pré-analyse des colonnes}
% ─────────────────────────────────────────────────────────────

\subsection{Objectif}

Analyse la structure des fichiers Excel en staging et produit des rapports
CSV sur les variations de colonnes d'une période à l'autre. Étape exploratoire,
sans modification des données.

\subsection{Variables configurables}

\begin{tabularx}{\textwidth}{lX}
  \toprule
  \textbf{Variable} & \textbf{Description} \\
  \midrule
  \texttt{BUCKET} & Bucket source. Défaut : \texttt{staging} \\
  \texttt{PREFIX\_MENSUELS} & Préfixe des fichiers mensuels.
    Défaut : \texttt{panel\_admin/fichiers\_mensuels} \\
  \texttt{PREFIX\_SORTIE} & Préfixe de sortie des rapports.
    Défaut : \texttt{panel\_admin/pre\_analyse} \\
  \bottomrule
\end{tabularx}

\subsection{Logique technique}

\textbf{Listage et tri.} Les fichiers \texttt{.xlsx} sont listés via
paginateur S3 et triés par nom (ordre alphabétique = ordre chronologique,
car les noms sont au format \texttt{MMYYYY}).

\textbf{Lecture des en-têtes.} Chaque fichier est téléchargé dans un fichier
temporaire (\texttt{download\_file} puis \texttt{NamedTemporaryFile}), puis lu
avec :

\begin{codebox}[]
f1 = pl.read\_excel(tmp\_path, sheet\_id=1, read\_options={"n\_rows": 0})
f2 = pl.read\_excel(tmp\_path, sheet\_id=2, read\_options={"n\_rows": 0})
\end{codebox}

Le paramètre \texttt{read\_options=\{"n\_rows": 0\}} (et non
\texttt{engine\_options}) charge uniquement les en-têtes sans les données,
pour des raisons de performance.

\textbf{Extraction des codes numériques F2.} Les colonnes dont le nom est
purement numérique sont identifiées par :

\begin{codebox}[]
codes\_num\_f2 = [c for c in f2.columns if str(c).strip().isdigit()]
\end{codebox}

\textbf{Gestion des erreurs.} Un fichier en erreur est ajouté à la liste
\texttt{erreurs} et n'arrête pas l'analyse des fichiers suivants.

\subsection{Analyses réalisées}

\begin{itemize}
  \item Colonnes apparues / disparues entre deux périodes consécutives
        (Feuille 1)
  \item Codes numériques apparus / disparus entre deux périodes
        consécutives (Feuille 2)
  \item Statistiques globales (min/max/moyenne du nombre de colonnes par
        fichier, codes F2 min/max)
\end{itemize}

\subsection{Fichiers produits}

\begin{tabularx}{\textwidth}{lX}
  \toprule
  \textbf{Fichier} & \textbf{Contenu} \\
  \midrule
  \texttt{colonnes\_f1\_uniques.csv} & Liste de toutes les colonnes F1
    distinctes sur l'ensemble des périodes \\
  \texttt{codes\_f2\_uniques.csv} & Liste de tous les codes numériques F2
    distincts \\
  \texttt{stats\_par\_periode.csv} & Nombre de colonnes F1/F2/codes par
    période \\
  \texttt{changements\_f1.csv} & Colonnes apparues / disparues F1
    (si changements détectés) \\
  \texttt{changements\_codes\_f2.csv} & Codes numériques apparus / disparus F2
    (si changements détectés) \\
  \bottomrule
\end{tabularx}

Tous ces fichiers sont déposés dans
\bucket{staging/panel\_admin/pre\_analyse/}.

% ─────────────────────────────────────────────────────────────
\section{02\_staging\_to\_bronze.py — Staging vers Bronze}
% ─────────────────────────────────────────────────────────────

\subsection{Objectif}

Convertit chaque fichier Excel mensuel en un fichier Parquet Bronze. Bronze =
données brutes, tout en \texttt{String} (\texttt{Utf8}), sans transformation
métier. Un fichier Parquet par période est produit.

\subsection{Variables configurables}

\begin{tabularx}{\textwidth}{lX}
  \toprule
  \textbf{Variable} & \textbf{Description} \\
  \midrule
  \texttt{BUCKET\_STAGING} & Bucket source. Défaut : \texttt{staging} \\
  \texttt{BUCKET\_BRONZE} & Bucket cible. Défaut : \texttt{bronze} \\
  \texttt{PREFIX\_MENSUELS} &
    Défaut : \texttt{panel\_admin/fichiers\_mensuels} \\
  \texttt{PREFIX\_BRONZE} & Défaut : \texttt{panel\_admin} \\
  \bottomrule
\end{tabularx}

\subsection{Fonctions utilitaires}

\subsubsection{normaliser\_nom\_colonne(nom)}

Normalise un nom de colonne brut en identifiant standardisé :
\begin{enumerate}[nosep]
  \item Supprime les préfixes du type \texttt{SHEET1.} via regex
        \texttt{\^{}[A-Z\_]+[0-9]+\textbackslash.}
  \item Passe en majuscules
  \item Remplace les caractères \texttt{-}, \texttt{/}, \texttt{\_},
        espace par \texttt{\_}
  \item Remplace les points par \texttt{\_}
  \item Normalisation Unicode NFKD + suppression des combinants
        (suppression des accents)
  \item Supprime tous les caractères non alphanumériques sauf \texttt{\_}
  \item Déduplique les \texttt{\_} consécutifs
  \item Supprime les \texttt{\_} en début/fin
\end{enumerate}

\subsubsection{deduplication\_noms(noms)}

Rend une liste de noms de colonnes unique en suffixant les doublons avec
\texttt{\_2}, \texttt{\_3}, etc. Exemple :
\texttt{[NOM, NOM, NOM]} $\rightarrow$ \texttt{[NOM, NOM\_2, NOM\_3]}.

\subsubsection{detecter\_entete(fichier, sheet\_id)}

Détecte automatiquement la ligne d'en-tête contenant ``matricule'' :
\begin{enumerate}[nosep]
  \item Lit les 20 premières lignes sans en-tête
        (\texttt{header\_row=None, n\_rows=20})
  \item Parcourt chaque ligne à la recherche du mot ``matricule''
        (insensible à la casse)
  \item Renvoie le numéro de ligne (0-indexé) si trouvé
  \item Défaut : ligne 1 pour Feuille 1, ligne 0 pour Feuille 2
\end{enumerate}

\subsection{Logique de traitement}

Pour chaque fichier Excel mensuel :

\begin{enumerate}
  \item Téléchargement dans un fichier temporaire (\texttt{download\_file})
  \item Détection de la ligne d'en-tête F1 et F2 via
        \texttt{detecter\_entete()}
  \item Lecture Feuille 1 avec \texttt{header\_row=skip\_f1}
  \item Normalisation et déduplication des noms de colonnes F1
  \item Lecture Feuille 2 avec \texttt{header\_row=skip\_f2}
  \item Jointure F2 sur la clé \textbf{matricule\textbar\textbar{}code\_organisme} :
    \begin{itemize}[nosep]
      \item La colonne clé F2 est détectée par regex
            \texttt{MATRICULE.*\textbackslash|\textbackslash|}
      \item Renommée en \texttt{CLE\_UNIQUE\_F2}
      \item Toutes les autres colonnes F2 reçoivent le suffixe
            \texttt{\_F2} (ex: \texttt{124\_F2}, \texttt{228\_F2})
      \item Jointure \texttt{left} sur \texttt{CLE\_UNIQUE} ou
            \texttt{MATRICULE} selon disponibilité
    \end{itemize}
  \item Ajout de deux colonnes de métadonnées :
    \begin{itemize}[nosep]
      \item \colname{PERIODE} : nom du fichier sans extension
            (format \texttt{MMYYYY})
      \item \colname{FICHIER\_SOURCE} : nom du fichier avec extension
    \end{itemize}
  \item Cast de toutes les colonnes en \texttt{Utf8} (String)
  \item Écriture Parquet avec compression Snappy :
        \bucket{bronze/panel\_admin/\{periode\}.parquet}
  \item Suppression du fichier temporaire
\end{enumerate}

\subsection{Nommage du fichier Parquet}

Le nom du fichier Parquet Bronze est le nom du fichier Excel source sans
l'extension, soit le format \texttt{MMYYYY} (ex: \texttt{012015.parquet}).

\subsection{Comportement en cas d'erreur}

Un fichier en erreur incrémente le compteur \texttt{nb\_err} et n'arrête
pas le traitement des autres fichiers. Le fichier temporaire est supprimé
même en cas d'erreur (bloc \texttt{finally}).

\subsection{Fichiers produits}

\bucket{bronze/panel\_admin/MMYYYY.parquet} — un fichier par mois.

% ─────────────────────────────────────────────────────────────
\section{03\_bronze\_to\_silver.py — Bronze vers Silver}
% ─────────────────────────────────────────────────────────────

\subsection{Objectif}

Applique la logique métier sur les données brutes Bronze : mapping des
colonnes vers des noms standardisés, normalisation des situations
administratives, enrichissement avec les codes ANSTAT, cast en types
numériques, suivi du taux de matching par période.

\subsection{Variables configurables}

\begin{tabularx}{\textwidth}{lX}
  \toprule
  \textbf{Variable} & \textbf{Description} \\
  \midrule
  \texttt{BUCKET\_STAGING} & \texttt{staging} \\
  \texttt{BUCKET\_BRONZE} & \texttt{bronze} \\
  \texttt{BUCKET\_SILVER} & \texttt{silver} \\
  \texttt{PREFIX\_BRONZE} & \texttt{panel\_admin} \\
  \texttt{PREFIX\_SILVER} & \texttt{panel\_admin} \\
  \texttt{PREFIX\_REFS} & \texttt{panel\_admin/references} \\
  \texttt{PREFIX\_VALID} & \texttt{panel\_admin/validation} \\
  \texttt{CODES\_A\_SURVEILLER} &
    \texttt{["CODE\_ORGANISME", "CODE\_GRADE", "CODE\_EMPLOI"]} ---
    codes dont le taux de matching est loggué en temps réel \\
  \bottomrule
\end{tabularx}

\subsection{Fonctions utilitaires}

\subsubsection{normaliser\_pour\_matching(texte)}

Normalisation textuelle utilisée pour le fuzzy join sur les libellés ANSTAT :
\begin{enumerate}[nosep]
  \item Remplace les tirets par des espaces
  \item Remplace les caractères \texttt{.,;:()\textbackslash[]/ } par des
        espaces
  \item Remplace les apostrophes (droite, courbe, unicode U+2019) par des
        espaces
  \item Passage en majuscules
  \item Normalisation Unicode NFKD + suppression des combinants
  \item Normalisation des espaces multiples
\end{enumerate}

\subsubsection{mapper\_colonnes(noms\_colonnes)}

Mappe chaque colonne brute vers un nom standardisé via 26 patrons regex.
Le tableau suivant liste le mapping complet :

\begin{longtable}{lp{7.5cm}}
  \toprule
  \textbf{Nom standardisé} & \textbf{Patterns regex} \\
  \midrule
  \endhead
  \texttt{cle\_unique} &
    \texttt{MATRICULE.*CODE.*ORGANISME}, \texttt{MATRICULE.*\textbar\textbar.*CODE} \\
  \texttt{matricule} & \texttt{\^{}MATRICULE\$}, \texttt{\^{}MATRIC\$} \\
  \texttt{nom} & \texttt{\^{}NOM\$}, \texttt{\^{}NOM\_PRENOM\$},
    \texttt{\^{}NOM\_ET\_PRENOM\$} \\
  \texttt{date\_naissance} & \texttt{DATE.*NAISSANCE}, \texttt{DATE\_NAISS},
    \texttt{NAISSANCE} \\
  \texttt{sexe} & \texttt{\^{}SEXE\$} \\
  \texttt{situation\_matrimoniale} & \texttt{SITUATION.*MATRIMONIALE},
    \texttt{STATUT.*MATRIMONIAL} \\
  \texttt{nombre\_enfant} & \texttt{NOMBRE.*ENFANT}, \texttt{NBR.*ENFANT} \\
  \texttt{situation} & \texttt{\^{}SITUATION\$},
    \texttt{SITUATION\_ADMINISTRATIVE} \\
  \texttt{date\_debut\_situation} & \texttt{DATE.*DEBUT.*SITUATION},
    \texttt{DATE\_SITUATION} \\
  \texttt{montant\_brut} & \texttt{MONTANT\_BRUT}, \texttt{SALAIRE\_BRUT},
    \texttt{REMUNERATION\_BRUTE}, \texttt{\^{}BRUT\$} \\
  \texttt{montant\_net} & \texttt{MONTANT\_NET}, \texttt{SALAIRE\_NET},
    \texttt{REMUNERATION\_NETTE}, \texttt{\^{}NET\$} \\
  \texttt{retenue\_pension} & \texttt{RETENUE.*PENSION}, \texttt{\^{}PENSION\$},
    \texttt{COTISATION.*PENSION} \\
  \texttt{impot} & \texttt{\^{}IMPOT\$}, \texttt{\^{}IGR\$},
    \texttt{IMPOT.*REVENU} \\
  \texttt{charge\_patronale} & \texttt{CHARGE.*PATRONALE},
    \texttt{CHARGES.*PATRONALES} \\
  \texttt{organisme} & \texttt{\^{}ORGANISME\$}, \texttt{MINISTERE},
    \texttt{DIRECTION} \\
  \texttt{lieu\_affectation} & \texttt{LIEU.*AFFECTATION},
    \texttt{\^{}AFFECTATION\$} \\
  \texttt{service} & \texttt{\^{}SERVICE\$} \\
  \texttt{emploi} & \texttt{\^{}EMPLOI\$}, \texttt{\^{}CORPS\$},
    \texttt{CORPS.*EMPLOI} \\
  \texttt{fonction} & \texttt{\^{}FONCTION\$} \\
  \texttt{grade} & \texttt{CLASSE.*ECHELON}, \texttt{CLASSE\_ECHELON} \\
  \texttt{statut\_fonctionnaire} & \texttt{\^{}GRADE\$} \\
  \texttt{poste} & \texttt{\^{}POSTE\$}, \texttt{LIBELLE.*POSTE} \\
  \texttt{prise\_service} & \texttt{PRISE.*SERVICE},
    \texttt{DATE.*PRISE.*SERVICE} \\
  \texttt{date\_retraite} & \texttt{DATE.*RETRAITE} \\
  \texttt{age\_retraite} & \texttt{AGE.*RETRAITE} \\
  \texttt{mois\_annee} & \texttt{MOIS.*ANNEE}, \texttt{\^{}PERIODE\$} \\
  \bottomrule
\end{longtable}

Si aucun patron ne correspond, le nom normalisé (minuscules) est conservé.

\subsubsection{normaliser\_situation(situation)}

Mappe la valeur brute de situation vers une catégorie normalisée :
\begin{itemize}[nosep]
  \item \texttt{en\_activite} : \texttt{EN ACTIVITE}, \texttt{ACTIVITE}
  \item \texttt{regul\_indemnites} : \texttt{REGUL. INDEMNITES},
        \texttt{REGUL INDEMNITES}, \texttt{REGULARISATION INDEMNITES}
  \item \texttt{demi\_solde} : \texttt{DEMI-SOLDE}, \texttt{DEMI SOLDE},
        \texttt{DEMISOLDE}, \texttt{1/2 SOLDE}
  \item \texttt{autre} : toute valeur non reconnue (y compris null, vide)
\end{itemize}

\subsubsection{charger\_table\_codes(...)}

Charge un onglet du fichier ANSTAT, identifie la colonne libellé par pattern,
la normalise pour le matching, et renvoie un dictionnaire avec la table,
le nom de la colonne code, et la colonne de jointure source.

\subsection{Chargement des tables ANSTAT}

Huit tables de référence sont chargées depuis le fichier ANSTAT :

\begin{tabularx}{\textwidth}{llX}
  \toprule
  \textbf{Code} & \textbf{Onglet} & \textbf{Colonnes extra} \\
  \midrule
  \texttt{CODE\_AFFECTATION} & \textit{lieu affectation} &
    CODE\_RÉGION, CODE\_SITE, CODE\_DRB, CODE\_DÉPARTEMENT,
    MONTANT TRANSPORT \\
  \texttt{CODE\_ORGANISME} & \textit{ORGANISME\_OK} &
    CODE\_ÉTABLISSEMENT, CODE\_ACCT \\
  \texttt{CODE\_POSITION\_SOLDE} & \textit{CODE\_SITUATION\_SOLDE} & --- \\
  \texttt{CODE\_EMPLOI} & \textit{HISTORIQUE\_ECHELLES\_CORPS} &
    GRADE\_ADMINISTRATIF\_ASSOCIE, Code\_CITP, Metier\_CITP \\
  \texttt{CODE\_SERVICE} & \textit{SERVICE} & --- \\
  \texttt{CODE\_FONCTION} & \textit{FONCTION} & --- \\
  \texttt{CODE\_GRADE} & \textit{GRADE} & --- \\
  \texttt{CODE\_POSTE} & \textit{LIBELLE POSTE} & --- \\
  \bottomrule
\end{tabularx}

\subsection{Logique de traitement par période}

Pour chaque fichier Bronze mensuel :

\begin{enumerate}
  \item Lecture du Parquet depuis MinIO Bronze
  \item Mapping des colonnes via \texttt{mapper\_colonnes()},
        avec déduplication des noms standards si collision
  \item \textbf{Forçage de mois\_annee} depuis le nom du fichier source
        (plus fiable que la valeur interne --- certains fichiers Bronze
        contiennent un \texttt{mois\_annee} erroné car il s'agit de doublons)
  \item Normalisation des situations et filtrage : les lignes
        \texttt{autre} sont \textbf{supprimées}
  \item Enrichissement ANSTAT : pour chaque code, jointure gauche sur
        libellé normalisé ; calcul du taux de matching
  \item Cast numérique des montants et des codes numériques (colonnes
        commençant par un chiffre)
  \item Accumulation dans \texttt{frames\_par\_annee}
\end{enumerate}

\subsection{Rapport de matching ANSTAT}

Un rapport \texttt{taux\_matching\_anstat\_par\_periode.csv} est sauvegardé
dans \bucket{staging/panel\_admin/validation/}. Il contient, pour chaque
période, le taux de matching (en \%) pour chacun des 8 codes ANSTAT.

Les codes \texttt{CODES\_A\_SURVEILLER} (CODE\_ORGANISME, CODE\_GRADE,
CODE\_EMPLOI) voient leur taux loggué en temps réel pendant l'exécution.

\subsection{Fichiers produits}

\bucket{silver/panel\_admin/YYYY.parquet} --- un fichier par année.

% ─────────────────────────────────────────────────────────────
\section{04\_validation\_silver.py — Validation qualité}
% ─────────────────────────────────────────────────────────────

\subsection{Objectif}

Contrôles qualité sur les tables Silver. Étape \textbf{optionnelle} dans
l'orchestrateur : un échec n'arrête pas le pipeline.

\subsection{Variables configurables}

\begin{tabularx}{\textwidth}{lX}
  \toprule
  \textbf{Variable} & \textbf{Description} \\
  \midrule
  \texttt{COLONNES\_CLES} &
    Colonnes dont la complétude est contrôlée toutes années confondues :
    \texttt{matricule}, \texttt{nom}, \texttt{montant\_brut},
    \texttt{montant\_net}, \texttt{CODE\_ORGANISME}, \texttt{CODE\_GRADE},
    \texttt{CODE\_EMPLOI}, \texttt{situation\_normalisee},
    \texttt{mois\_annee} \\
  \texttt{COLONNES\_NA\_SURVEY} &
    Colonnes dont le taux de NA est suivi par année :
    \texttt{matricule}, \texttt{nom}, \texttt{montant\_brut},
    \texttt{montant\_net}, \texttt{organisme}, \texttt{grade},
    \texttt{emploi}, \texttt{lieu\_affectation}, \texttt{service} \\
  \bottomrule
\end{tabularx}

\subsection{Contrôles réalisés}

\begin{enumerate}
  \item \textbf{Complétude des colonnes clés} (toutes années confondues).
        Alerte si taux de complétion $<$ 95\%.

  \item \textbf{Colonnes suspectes} (NA variable selon l'année).
        Seuil : taux NA $>$ 95\% au max ET $<$ 5\% au min (variation $>$ 90
        points entre années). Si des colonnes suspectes sont détectées,
        le rapport \texttt{na\_par\_annee\_colonnes\_surveillees.csv} est
        sauvegardé.

  \item \textbf{Doublons matricule $\times$ mois\_annee}. Le rapport
        \texttt{doublons\_matricule\_periode.csv} est sauvegardé si des
        doublons existent.

  \item \textbf{Cohérence des montants} (net $>$ brut) : comptage des lignes
        incohérentes.

  \item \textbf{Distribution des situations administratives} : répartition
        \texttt{en\_activite} / \texttt{regul\_indemnites} /
        \texttt{demi\_solde}.
\end{enumerate}

\subsection{Résumé final}

À l'issue de la validation, le script affiche le nombre d'agents uniques
(matricules distincts), la plage temporelle, et le nombre de périodes
distinctes présentes dans le Silver.

\subsection{Fichiers produits (si anomalies)}

\begin{itemize}[nosep]
  \item \bucket{staging/panel\_admin/validation/na\_par\_annee\_colonnes\_surveillees.csv}
  \item \bucket{staging/panel\_admin/validation/doublons\_matricule\_periode.csv}
\end{itemize}

% ─────────────────────────────────────────────────────────────
\section{05\_silver\_to\_gold.py — Silver vers Gold}
% ─────────────────────────────────────────────────────────────

\subsection{Objectif}

Produit deux tables Gold d'agrégats à partir des tables Silver annuelles.

\subsection{Variables configurables}

\begin{tabularx}{\textwidth}{lX}
  \toprule
  \textbf{Variable} & \textbf{Description} \\
  \midrule
  \texttt{BUCKET\_SILVER} & \texttt{silver} \\
  \texttt{BUCKET\_GOLD} & \texttt{gold} \\
  \texttt{BUCKET\_STAGING} & \texttt{staging} \\
  \texttt{KEY\_GOLD\_MS} &
    \texttt{panel\_admin/masse\_salariale.parquet} \\
  \texttt{KEY\_GOLD\_EFF} &
    \texttt{panel\_admin/effectifs.parquet} \\
  \texttt{PREFIX\_EXPORTS} &
    \texttt{panel\_admin/exports\_gold} \\
  \bottomrule
\end{tabularx}

\subsection{Table Gold 1 : masse\_salariale.parquet}

Agrégation par \texttt{mois\_annee} $\times$ \texttt{annee} $\times$
\texttt{mois} $\times$ \texttt{code\_organisme} $\times$
\texttt{organisme} $\times$ \texttt{situation\_normalisee}.

Colonnes produites :

\begin{tabularx}{\textwidth}{lX}
  \toprule
  \textbf{Colonne} & \textbf{Définition} \\
  \midrule
  \texttt{nb\_lignes} & Nombre de lignes dans le groupe \\
  \texttt{nb\_agents\_uniques} & Nombre de matricules distincts \\
  \texttt{masse\_brute} & Somme des montants bruts \\
  \texttt{masse\_nette} & Somme des montants nets \\
  \texttt{salaire\_brut\_moyen} & Moyenne du montant brut \\
  \texttt{salaire\_net\_moyen} & Moyenne du montant net \\
  \texttt{salaire\_brut\_mediane} & Médiane du montant brut (P50) \\
  \texttt{salaire\_net\_mediane} & Médiane du montant net (P50) \\
  \bottomrule
\end{tabularx}

\textit{Note :} seules les lignes où \texttt{montant\_brut} ET
\texttt{montant\_net} sont non nuls participent à cette agrégation.
\texttt{code\_organisme} vaut \texttt{NON\_CODE} si null, \texttt{organisme}
vaut \texttt{NON\_IDENTIFIE} si null.

\subsection{Table Gold 2 : effectifs.parquet}

Agrégation par \texttt{mois\_annee} $\times$ \texttt{annee} $\times$
\texttt{mois} $\times$ \texttt{situation\_normalisee} $\times$
\texttt{code\_organisme}.

Colonnes produites :

\begin{tabularx}{\textwidth}{lX}
  \toprule
  \textbf{Colonne} & \textbf{Définition} \\
  \midrule
  \texttt{nb\_lignes} & Nombre de lignes dans le groupe \\
  \texttt{nb\_agents\_uniques} & Nombre de matricules distincts \\
  \texttt{nb\_hommes} & Nombre de lignes où \texttt{sexe == "1"} \\
  \texttt{nb\_femmes} & Nombre de lignes où \texttt{sexe == "2"} \\
  \bottomrule
\end{tabularx}

\subsection{Export CSV}

Deux exports CSV sont déposés dans
\bucket{staging/panel\_admin/exports\_gold/} :
\begin{itemize}[nosep]
  \item \texttt{masse\_salariale\_par\_periode\_organisme.csv}
  \item \texttt{effectifs\_par\_periode.csv}
\end{itemize}

% ─────────────────────────────────────────────────────────────
\section{06\_compiler\_panel.py — Compilation du panel}
% ─────────────────────────────────────────────────────────────

\subsection{Objectif}

Construit le panel individu $\times$ période à partir des tables Silver
annuelles. Produit des fichiers annuels et un fichier complet en streaming.

\subsection{Colonnes du panel (ordre exact)}

\begin{longtable}{lll}
  \toprule
  \textbf{Groupe} & \textbf{Colonne} & \textbf{Type} \\
  \midrule
  \endhead
  Identité & \texttt{matricule} & Utf8 \\
  & \texttt{nom} & Utf8 \\
  & \texttt{sexe} & Utf8 \\
  & \texttt{date\_naissance} & Utf8 \\
  & \texttt{situation\_matrimoniale} & Utf8 \\
  & \texttt{nombre\_enfant} & Utf8 \\
  \midrule
  Position & \texttt{situation\_normalisee} & Utf8 \\
  & \texttt{CODE\_ORGANISME} & Utf8 \\
  & \texttt{organisme} & Utf8 \\
  & \texttt{CODE\_EMPLOI} & Utf8 \\
  & \texttt{emploi} & Utf8 \\
  & \texttt{Code\_CITP} & Utf8 \\
  & \texttt{Metier\_CITP} & Utf8 \\
  & \texttt{CODE\_GRADE} & Utf8 \\
  & \texttt{grade} & Utf8 \\
  & \texttt{statut\_fonctionnaire} & Utf8 \\
  & \texttt{CODE\_FONCTION} & Utf8 \\
  & \texttt{fonction} & Utf8 \\
  & \texttt{CODE\_SERVICE} & Utf8 \\
  & \texttt{service} & Utf8 \\
  & \texttt{CODE\_AFFECTATION} & Utf8 \\
  & \texttt{lieu\_affectation} & Utf8 \\
  & \texttt{CODE\_POSTE} & Utf8 \\
  & \texttt{prise\_service} & Utf8 \\
  & \texttt{date\_retraite} & Utf8 \\
  & \texttt{age\_retraite} & Utf8 \\
  \midrule
  Temporel & \texttt{mois\_annee} & String (brut Silver) \\
  & \texttt{annee} & Int32 \\
  & \texttt{mois} & Int32 \\
  \midrule
  Salaires & \texttt{montant\_brut} & Float64 \\
  & \texttt{montant\_net} & Float64 \\
  & \texttt{retenue\_pension} & Float64 \\
  & \texttt{impot} & Float64 \\
  & \texttt{charge\_patronale} & Float64 \\
  \midrule
  Traçabilité & \texttt{fichier\_source} & Utf8 \\
  \bottomrule
\end{longtable}

\begin{infobox}[Note sur GRADE et GRADE\_SOURCE]
  Ces deux colonnes sont ajoutées par l'étape \textbf{06b}
  (\texttt{06b\_imputer\_grades.py}), pas par 06. Elles sont placées juste
  après \texttt{statut\_fonctionnaire} dans l'ordre final.
\end{infobox}

\subsection{Fonction construire\_panel\_annee(df)}

Pour chaque table Silver annuelle :

\begin{enumerate}
  \item Extraction de \texttt{annee} et \texttt{mois} depuis
        \texttt{mois\_annee}. Deux formats supportés :
    \begin{itemize}[nosep]
      \item Format MMYYYY (6 caractères) : \texttt{annee = slice(2,4)},
            \texttt{mois = slice(0,2)}
      \item Format YYYY-MM : \texttt{annee = slice(0,4)},
            \texttt{mois = slice(5,2)}
    \end{itemize}
  \item Sélection des colonnes présentes dans \texttt{COLONNES\_PANEL}
        (ordre préservé, colonnes absentes ignorées)
  \item Déduplication sur \texttt{matricule $\times$ mois\_annee}
        (garde la \textbf{première} occurrence)
  \item Tri par \texttt{matricule} $\rightarrow$ \texttt{annee}
        $\rightarrow$ \texttt{mois}
\end{enumerate}

\subsection{Fusion complète en streaming}

Le fichier \texttt{panel\_complet.parquet} est produit par :
\begin{enumerate}
  \item Téléchargement de tous les fichiers \texttt{panel\_YYYY.parquet}
        dans un répertoire temporaire (\texttt{TemporaryDirectory})
  \item \texttt{pl.scan\_parquet(chemins\_locaux)} : lecture en mode lazy,
        sans charger les données en mémoire
  \item \texttt{.sort([...]).sink\_parquet(chemin, compression="snappy")} :
        écriture en streaming, jamais tout en RAM
  \item Upload du fichier local vers MinIO Gold
\end{enumerate}

Cette approche permet de traiter le panel complet (millions de lignes) sans
surcharger la mémoire.

\subsection{Statistiques affichées}

Par année : \texttt{nb\_obs}, \texttt{nb\_individus}, \texttt{brut\_moyen},
\texttt{brut\_mediane}.

\subsection{Fichiers produits}

\begin{itemize}[nosep]
  \item \bucket{gold/panel\_admin/panel\_YYYY.parquet} (un par année)
  \item \bucket{gold/panel\_admin/panel\_complet.parquet}
\end{itemize}

% ─────────────────────────────────────────────────────────────
\section{06b\_imputer\_grades.py — Imputation des grades}
% ─────────────────────────────────────────────────────────────

\subsection{Objectif}

Calcule les grades manquants ou invalides par imputation statistique en
cascade sur la colonne \texttt{emploi}, puis met à jour tous les fichiers
Gold (panel annuels + panel complet).

\subsection{Pourquoi une étape dédiée ?}

L'imputation de niveau P3 (global, toutes années) est d'autant plus fiable
qu'elle s'appuie sur l'intégralité de l'historique. Travailler année par
année donnerait des lookups instables pour les emplois rares.

\subsection{Constantes}

\begin{codebox}[]
GRADES\_VALIDES = A1..A7, B1..B6, C1..C5, D1..D3  # 21 grades
GRADES\_INVALIDES = ["B7", "C7", "D7", "A", "B", "C", "D"]
\end{codebox}

\subsection{Fonction extraire\_grade(statut)}

Extrait le grade depuis \texttt{statut\_fonctionnaire} via l'expression
régulière \texttt{Cat[eé]gorie\textbackslash{}s+([ABCD][0-9]*)}
(insensible à la casse). Exemples : ``Catégorie A3'' $\rightarrow$
\texttt{A3}, ``Catégorie b2'' $\rightarrow$ \texttt{B2}.

Si aucun grade n'est trouvé, retourne \texttt{"NF"}.

\textit{Note :} les grades NF (non-fonctionnaires ou valeur non reconnue)
restent \texttt{NF} et ne sont pas imputés.

\subsection{Séquence en 6 étapes}

\textbf{Étape 1.} Chargement du \texttt{panel\_complet.parquet}.

\textbf{Étape 2.} Extraction du grade brut via \texttt{extraire\_grade()}.
Bilan : grades valides / invalides / NF affiché.

\textbf{Étape 3.} Imputation en cascade (si \texttt{n\_invalides > 0}) :
\begin{itemize}
  \item Extraction de \texttt{\_ANNEE} (slice 2:4 de mois\_annee) et
        \texttt{\_MOIS\_NUM} (slice 0:2) pour les lookups
  \item Construction de la base des grades valides uniquement
  \item \textbf{Lookup P1} : \texttt{GRADE.mode().first()} groupé par
        \texttt{emploi $\times$ \_ANNEE $\times$ \_MOIS\_NUM}
  \item \textbf{Lookup P2} : \texttt{GRADE.mode().first()} groupé par
        \texttt{emploi $\times$ \_ANNEE}
  \item \textbf{Lookup P3} : \texttt{GRADE.mode().first()} groupé par
        \texttt{emploi} (global, toutes années)
  \item Application de la cascade : si grade invalide, tenter P1 $\rightarrow$
        P2 $\rightarrow$ P3, conserver l'invalide si aucun lookup n'aboutit
\end{itemize}

\textbf{Étape 4.} Réordonnancement : \colname{GRADE} et
\colname{GRADE\_SOURCE} placés juste après \texttt{statut\_fonctionnaire}.

Colonnes temporaires supprimées : \texttt{\_ANNEE}, \texttt{\_MOIS\_NUM},
\texttt{GRADE\_P1}, \texttt{GRADE\_P2}, \texttt{GRADE\_P3},
\texttt{GRADE\_SOURCE\_TMP}.

\textbf{Étape 5.} Mise à jour des fichiers Gold annuels : jointure sur
\texttt{(matricule, mois\_annee)} pour injecter \colname{GRADE} et
\colname{GRADE\_SOURCE}.

\textbf{Étape 6.} Mise à jour du \texttt{panel\_complet.parquet}.

\subsection{Valeurs de GRADE\_SOURCE}

\begin{tabularx}{\textwidth}{lX}
  \toprule
  \textbf{Valeur} & \textbf{Signification} \\
  \midrule
  \texttt{OBSERVE} & Grade extrait directement de \texttt{statut\_fonctionnaire},
    valide \\
  \texttt{IMPUTE} & Grade remplacé par le mode statistique de l'emploi (P1/P2/P3) \\
  \texttt{NF\_NON\_IMPUTABLE} & Aucun grade valide trouvé ni imputé
    (non-fonctionnaire ou emploi sans historique de grade valide) \\
  \bottomrule
\end{tabularx}

\subsection{Bilan affiché}

Pourcentage OBSERVE / IMPUTE / NF\_NON\_IMPUTABLE + nombre de grades encore
invalides après imputation.

% ─────────────────────────────────────────────────────────────
\section{07\_calcul\_indicateur.py — Calcul des indicateurs}
% ─────────────────────────────────────────────────────────────

\subsection{Objectif}

Produit les indicateurs statistiques complets sur les rémunérations à partir
des fichiers Gold panel. Requiert que \texttt{06b\_imputer\_grades.py} ait
été exécuté au préalable (colonnes GRADE et GRADE\_SOURCE obligatoires).

\subsection{Variables configurables}

\begin{tabularx}{\textwidth}{lX}
  \toprule
  \textbf{Variable} & \textbf{Défaut} & \\
  \midrule
  \texttt{MODE\_SALAIRE\_BRUT} & \texttt{True} &
    \texttt{True} = agrégation par \texttt{cle\_unique} (salaire brut) ;
    \texttt{False} = agrégation par \texttt{matricule} (revenu salarial) \\
  \texttt{PERIODE\_RECENTE} & \texttt{True} &
    \texttt{True} = années 2024--2025 ; \texttt{False} = 2015 $\rightarrow$
    dernière année \\
  \texttt{INCLURE\_ZEROS} & \texttt{True} &
    \texttt{True} = inclure salaires = 0 ou NA \\
  \texttt{SEUIL\_WINS} & \texttt{0.01} &
    Seuil de winsorisation P1/P99. Mettre \texttt{None} pour désactiver \\
  \texttt{MIN\_OBS\_WINS} & \texttt{100} &
    Effectif minimum par cellule mois $\times$ grade $\times$ CITP pour
    appliquer la winsorisation CITP \\
  \texttt{KEY\_CORRESPONDANCE} &
    \texttt{panel\_admin/references/CORRESPONDANCE\_GRILLE\_EMPLOI.xlsx} &
    Clé S3 de la table barème \\
  \bottomrule
\end{tabularx}

\subsection{Composantes du salaire brut}

\begin{codebox}[]
COMPOSANTES\_SALAIRE = ["124\_f2", "228\_f2", "271\_f2", "421\_f2", "425\_f2"]
\end{codebox}

Ces colonnes correspondent aux codes de rémunération de la Feuille 2, avec
le suffixe \texttt{\_f2} ajouté lors de l'étape 02. Si aucune composante
n'est trouvée dans le Gold, le script utilise \texttt{montant\_brut} comme
fallback.

\subsection{Nom du fichier de sortie (dynamique)}

\begin{codebox}[]
indicateurs\_{mode\_label}\_{periode\_label}\_{zeros\_label}.xlsx
# Exemples :
indicateurs\_salaire\_brut\_2024\_2025\_avec\_zeros.xlsx
indicateurs\_revenu\_salarial\_2015\_recent\_sans\_zeros.xlsx
\end{codebox}

\subsection{Constantes barème}

\begin{codebox}[]
GRADES\_PRIMAIRE = \{"B3", "C3", "D1", "D2"\}

ORDRE\_MODELE\_ANSTAT = [
    ("TOUS SECTEURS CONFONDUS",                       "\_GLOBAL\_"),
    ("SECTEUR EDUCATION",                             None),
    ("Primaire",                                      "ENS SEC \& PRIM\_PRIM"),
    ("Secondaire",                                    "ENS SEC \& PRIM\_SEC"),
    ("Supérieur",                                     "ENS SUP"),
    ("SECTEUR SANTE",                                 None),
    ("Cadres supérieurs",  "Cadre-Sup-Santé Generaliste \& Specialiste"),
    ("Personnels Techniques",                         "Tech-Santé"),
    ("AUTRES SECTEURS",                               None),
    ("Para-militaire / Douanes / Eaux \& Forêts",     "police"),
    ("Barème Général",                                "Barême Général"),
    ("PERSONNELS SOUS STATUT PARTICULIER",            None),
    ("Magistrat",                                     "MAGISTRAT "),
    ("Corps préfectoral",                             "CORPS PREF"),
    ("Corps diplomatique",                            "CORPS DIPL"),
    ("Greffier",                                      "GREFFE"),
]
\end{codebox}

\subsection{Séquence complète en 12 étapes}

\subsubsection{Étape 1 — Chargement}

Chargement des fichiers \texttt{panel\_YYYY.parquet} (pas le
\texttt{panel\_complet}, pour éviter de le charger deux fois). Concaténation
diagonale en mémoire.

\subsubsection{Étape 2 — Vérification GRADE / GRADE\_SOURCE}

Si l'une des deux colonnes est absente, le script lève une
\texttt{RuntimeError} et demande d'exécuter \texttt{06b} d'abord.

\subsubsection{Étape 3 — Extraction ANNEE / MOIS}

Depuis \texttt{mois\_annee} en format MMYYYY :
\texttt{ANNEE = slice(2,4)}, \texttt{MOIS\_NUM = slice(0,2)}.
Filtre sur \texttt{ANNEES\_ANALYSE}.

\subsubsection{Étape 4 — Calcul SALAIRE\_BRUT}

Somme horizontale des composantes présentes. Fallback :
\texttt{montant\_brut}. Fallback ultime : 0.

\subsubsection{Étape 5 — Standardisation SEXE\_STD}

\begin{tabularx}{\textwidth}{lX}
  \toprule
  \textbf{SEXE\_STD} & \textbf{Valeurs reconnues} \\
  \midrule
  \texttt{Homme} & MASCULIN, M, H, HOMME, MALE, 1 \\
  \texttt{Femme} & FEMININ, F, FEMME, FEMALE, 2, FÉMININ, FEMININE \\
  \texttt{Inconnu} & Tout autre $\rightarrow$ \textbf{exclu des calculs} \\
  \bottomrule
\end{tabularx}

\subsubsection{Étape 6 — Détection multi-postes}

Comptage du nombre de lignes par \texttt{matricule $\times$ mois\_annee}
$\rightarrow$ colonne \texttt{NB\_POSTES}.

\subsubsection{Étape 7 — Filtrage grades valides}

Filtre sur \texttt{GRADE.is\_in(GRADES\_VALIDES)}.
Si \texttt{INCLURE\_ZEROS=False}, filtre additionnel \texttt{SALAIRE\_BRUT > 0}.
Les \texttt{NF\_NON\_IMPUTABLE} sont exclus.

\subsubsection{Étape 7b — Fuzzy join emploi $\rightarrow$ barème}

\textbf{(Nouveau)} Trois fonctions :
\begin{itemize}
  \item \texttt{charger\_table\_bareme()} : télécharge
        \texttt{CORRESPONDANCE\_GRILLE\_EMPLOI.xlsx} depuis staging, normalise
        \texttt{emploi\_ref} pour le matching, retourne une table
        \texttt{(emploi\_ref\_norm, bareme, grade\_ref)}. Si le fichier est
        absent : retourne \texttt{None}, \texttt{BAREME\_DISPONIBLE=False},
        feuilles barème ignorées.
  \item \texttt{fuzzy\_join\_bareme(df, table\_join, col\_emploi)} : ajoute les
        colonnes \texttt{bareme} et \texttt{grade\_ref} via jointure gauche sur
        libellé normalisé. Avertissement si taux $<$ 20\%.
  \item \texttt{ajouter\_bareme\_detail(df, col\_grade)} : scinde
        \texttt{ENS SEC \& PRIM} en \texttt{ENS SEC \& PRIM\_PRIM} (grades B/C/D
        dans \texttt{GRADES\_PRIMAIRE}) et \texttt{ENS SEC \& PRIM\_SEC} (grades A).
\end{itemize}

\subsubsection{Étape 8 — Agrégation REVENU\_TOTAL}

Agrégation par \texttt{(ANNEE, MOIS\_NUM, MOIS, cle\_unique|matricule)} :
\begin{itemize}[nosep]
  \item \texttt{REVENU\_TOTAL} = somme de \texttt{SALAIRE\_BRUT}
  \item \texttt{GRADE\_PRINCIPAL} : grade du poste le plus élevé (valeur
        ordinale : A=80+n, B=70+n, C=60+n, D=50+n)
  \item \texttt{GRADE\_SOURCE\_PRINCIPAL} : source du grade principal
  \item \texttt{SEXE\_STD}, \texttt{Code\_CITP}, \texttt{Metier\_CITP},
        \texttt{bareme}, \texttt{bareme\_detail}
\end{itemize}

\subsubsection{Étape 8b — Winsorisation REVENU\_TOTAL}

\textbf{(Nouveau)} Appliquée en mémoire uniquement (le Gold n'est pas modifié) :
\begin{itemize}
  \item Si n\_obs dans le groupe \texttt{mois $\times$ grade $\times$ CITP}
        $\geq$ \texttt{MIN\_OBS\_WINS} : clip à P1/P99 du groupe CITP
  \item Sinon : clip à P1/P99 du groupe \texttt{mois $\times$ grade} (repli)
  \item Crée la colonne \texttt{REVENU\_TOTAL\_W} utilisée pour les indicateurs
\end{itemize}

\subsubsection{Étape 9 — Indicateurs REVENU TOTAL}

Fonction \texttt{compute\_indicators()} : effectif, moyenne, médiane, min,
max, P25, P75 par groupe. Groupements produits :
\begin{itemize}[nosep]
  \item Par grade
  \item Par grade $\times$ sexe
  \item Par CITP $\times$ métier
  \item Par barème (si \texttt{BAREME\_DISPONIBLE})
  \item Par barème $\times$ grade (si \texttt{BAREME\_DISPONIBLE})
  \item Par barème $\times$ sexe (si \texttt{BAREME\_DISPONIBLE})
\end{itemize}

\subsubsection{Étape 10 — Indicateurs SALAIRE BRUT (pivot grade)}

\textbf{(Nouveau)} Calcul sur les lignes individuelles (non agrégées par agent) :
\begin{itemize}[nosep]
  \item Winsorisation identique à 8b mais sur \texttt{SALAIRE\_BRUT}
  \item Moyenne par \texttt{(ANNEE, MOIS\_NUM, MOIS, GRADE)} $\rightarrow$
        pivot grade $\rightarrow$ \texttt{brut\_moyen\_wide}
  \item Médiane par \texttt{(ANNEE, MOIS\_NUM, MOIS, GRADE)} $\rightarrow$
        pivot grade $\rightarrow$ \texttt{brut\_median\_wide}
\end{itemize}

\subsubsection{Étape 11 — Multi-postes}

Effectifs par \texttt{GRADE\_PRINCIPAL $\times$ NB\_POSTES} (pivot sur
\texttt{1/2/3/4/5+}). Distribution globale mono/multi-postes avec revenu moyen.

\subsubsection{Étape 12 — Export Excel}

Toutes les tables sont écrites dans un fichier Excel multi-feuilles avec
style d'en-tête bleu et gel de la première ligne.

\subsection{Feuilles produites}

\begin{longtable}{lX}
  \toprule
  \textbf{Feuille} & \textbf{Contenu} \\
  \midrule
  \endhead
  \texttt{REVENU\_Grade\_Detail} &
    Effectif, moy, méd, min, max, P25, P75 par
    \texttt{ANNEE $\times$ MOIS $\times$ GRADE\_PRINCIPAL} \\
  \texttt{REVENU\_Grade\_Sexe} &
    Idem par \texttt{GRADE $\times$ SEXE\_STD} \\
  \texttt{REVENU\_CITP\_Detail} &
    Idem par \texttt{Code\_CITP $\times$ Metier\_CITP}
    (si CITP disponible) \\
  \texttt{SALAIRE\_BRUT\_Moyen} &
    Salaire brut moyen par \texttt{ANNEE $\times$ MOIS} (pivot grade) \\
  \texttt{SALAIRE\_BRUT\_Median} &
    Salaire brut médian par \texttt{ANNEE $\times$ MOIS} (pivot grade) \\
  \texttt{MULTI\_Par\_Grade\_NbPostes} &
    Effectifs par \texttt{GRADE $\times$ NB\_POSTES}
    (pivot 1/2/3/4/5+) \\
  \texttt{MULTI\_Distribution} &
    Distribution globale mono/multi-postes avec revenu moyen \\
  \texttt{REVENU\_Bareme\_Detail} &
    Indicateurs par barème / grille salariale
    (si \texttt{BAREME\_DISPONIBLE}) \\
  \texttt{REVENU\_Bareme\_Grade} &
    Indicateurs par barème $\times$ grade
    (si \texttt{BAREME\_DISPONIBLE}) \\
  \texttt{REVENU\_Bareme\_Sexe} &
    Indicateurs par barème $\times$ sexe
    (si \texttt{BAREME\_DISPONIBLE}) \\
  \texttt{MODELE\_ANSTAT} &
    Format hiérarchique min/moyen/max par secteur (si barème disponible) \\
  \bottomrule
\end{longtable}

\subsection{Feuille MODELE\_ANSTAT}

Format calqué sur le modèle statistique de l'ANSTAT :
\begin{itemize}
  \item \textbf{Ligne 1} : titre fusionné A1:D1, fond bleu foncé
        (\texttt{1F4E79}), ``MODÈLE DE DONNÉES DE SALAIRE DES
        FONCTIONNAIRES POUR \{annee\}''
  \item \textbf{Ligne 2} : en-têtes (fond bleu \texttt{2E75B6}) :
        salaire brut minimum / moyen / maximum
  \item \textbf{TOUS SECTEURS CONFONDUS} : fond jaune \texttt{FFF2CC},
        valeurs globales
  \item \textbf{Titres secteur} (SECTEUR EDUCATION, SECTEUR SANTE, etc.) :
        fond bleu clair \texttt{D6E4F0}, cellule fusionnée A:D
  \item \textbf{Sous-groupes} : alternance blanc/bleu très clair
        \texttt{FFFFFF}/\texttt{EBF3FB}, indentation, format \texttt{\#,\#\#0}
  \item Largeurs : colonne A = 55, colonnes B-D = 22
\end{itemize}

Basé sur \texttt{SALAIRE\_BRUT} ligne par ligne (pas sur le
\texttt{REVENU\_TOTAL} agrégé par agent).

% ─────────────────────────────────────────────────────────────
\section{08\_creation\_fichier\_excel\_avec\_sommaire.py}
% ─────────────────────────────────────────────────────────────

\subsection{Objectif}

Lit le fichier d'indicateurs produit par 07 et génère un fichier Excel
navigable avec un sommaire à 5 sections et des feuilles de détail.

\begin{warningbox}[Synchronisation des paramètres]
  Les trois paramètres \texttt{MODE\_SALAIRE\_BRUT}, \texttt{PERIODE\_RECENTE}
  et \texttt{INCLURE\_ZEROS} doivent être \textbf{identiques} à ceux utilisés
  lors de l'exécution de \texttt{07\_calcul\_indicateur.py}, afin que le script
  retrouve le bon fichier source.
\end{warningbox}

\subsection{Fichiers source et sortie}

\begin{itemize}[nosep]
  \item \textbf{Source} :
        \bucket{staging/panel\_admin/exports\_gold/indicateurs\_\{mode\}\_\{periode\}\_\{zeros\}.xlsx}
  \item \textbf{Sortie} :
        \bucket{staging/panel\_admin/exports\_gold/indicateurs\_\{mode\}\_\{periode\}\_\{zeros\}\_SOMMAIRE.xlsx}
\end{itemize}

\subsection{Fonctions utilitaires Excel}

\begin{tabularx}{\textwidth}{lX}
  \toprule
  \textbf{Fonction} & \textbf{Rôle} \\
  \midrule
  \texttt{style\_cell(cell, fill, font, align)} &
    Applique fill, font et alignment en une seule ligne \\
  \texttt{merge\_titre(ws, texte, row, n\_cols, fill, font)} &
    Crée un titre fusionné stylisé sur \texttt{n\_cols} colonnes \\
  \texttt{lien\_interne(ws, row, col, texte, feuille\_cible)} &
    Crée un hyperlien Excel interne :
    \texttt{=HYPERLINK("\#FEUILLE!A1","texte")} \\
  \texttt{ecrire\_entetes(ws, entetes, row)} &
    Ligne d'en-tête stylisée (fond bleu, police blanche) \\
  \texttt{ecrire\_df\_dans\_feuille(ws, df, start\_row=4)} &
    Écrit un DataFrame à partir de la ligne 4 (lignes 1-3 réservées :
    titre, lien retour, vide) \\
  \texttt{ajouter\_lien\_retour(ws, row=2)} &
    Lien \texttt{<- Retour au sommaire} en ligne 2 \\
  \texttt{ajuster\_largeurs(ws, \{col\_index: width\})} &
    Définit les largeurs de colonnes par index \\
  \texttt{extract\_grand\_groupe(code\_citp)} &
    Extrait le premier chiffre d'un code CITP (grand groupe 1-9) \\
  \bottomrule
\end{tabularx}

\subsection{Structure du sommaire (5 sections)}

\textbf{Section 1 — CITP.}
En-têtes : Code, Grand Groupe CITP, Effectif total, Nb métiers.
9 grands groupes (1 = Directeurs...9 = Professions élémentaires).
Lien vers feuille \texttt{CITP\_GG\{n\}}.

\textbf{Section 2 — Grade.}
En-têtes : Grade, Description (Catégorie A/B/C/D), Effectif total,
Revenu moyen (pondéré). Lien vers feuille \texttt{GRADE\_\{grade\}}.

\textbf{Section 3 — Sexe.}
En-têtes : Sexe, Description, Effectif total, Revenu moyen (pondéré).
Lien vers feuille \texttt{SEXE\_\{Homme|Femme\}}.

\textbf{Section 4 — Multi-postes.}
En-têtes : Grade, Description, Mono-postes (NB\_POSTES=1),
Multi-postes (NB\_POSTES$\geq$2). Lien vers feuille
\texttt{MULTI\_\{grade\}}.

\textbf{Section 5 — Barème (nouveau, si \texttt{BAREME\_DISPONIBLE}).}
En-têtes : Barème / Grille salariale, Effectif total, Revenu moyen,
Détail grade (lien), Détail sexe (lien).
En bas de section : lien vers \texttt{MODELE\_ANSTAT}.

\subsection{Feuilles de détail créées}

\begin{tabularx}{\textwidth}{lX}
  \toprule
  \textbf{Feuilles} & \textbf{Contenu} \\
  \midrule
  \texttt{CITP\_GG1} \ldots \texttt{CITP\_GG9} &
    Données CITP filtrées par grand groupe \\
  \texttt{GRADE\_A1} \ldots \texttt{GRADE\_D3} &
    Indicateurs filtrés par grade \\
  \texttt{SEXE\_Homme}, \texttt{SEXE\_Femme} &
    Indicateurs filtrés par sexe \\
  \texttt{MULTI\_A1} \ldots \texttt{MULTI\_D3} &
    Multi-postes filtrés par grade \\
  \texttt{BAREME\_\{nom\}} &
    Indicateurs par barème (1 feuille par barème) \\
  \texttt{BAREME\_G\_\{nom\}} &
    Indicateurs barème $\times$ grade \\
  \texttt{BAREME\_S\_\{nom\}} &
    Indicateurs barème $\times$ sexe \\
  \texttt{MODELE\_ANSTAT} &
    Copie directe depuis le fichier 07 \\
  \bottomrule
\end{tabularx}

Chaque feuille de détail comporte en ligne 1 un titre fusionné, en ligne 2
un lien ``$\leftarrow$ Retour au sommaire'', et les données à partir de la
ligne 4.

% ═════════════════════════════════════════════════════════════
% CHAPITRE 5 — L'ORCHESTRATEUR
% ═════════════════════════════════════════════════════════════
\chapter{L'orchestrateur}

\section{Présentation}

\scriptname{orchestrateur.py} lance l'ensemble du pipeline ou une sélection
d'étapes. Il gère l'ordre d'exécution, la durée, et l'arrêt en cas d'échec.

\begin{warningbox}[Étape 00 non incluse dans l'orchestrateur]
  L'étape \texttt{00\_upload\_staging.py} est intentionnellement exclue de
  l'orchestrateur car elle ne doit être exécutée qu'une seule fois
  (dépôt initial). Lancez-la manuellement avant d'appeler l'orchestrateur.
\end{warningbox}

\section{Étapes gérées}

\begin{tabularx}{\textwidth}{llX}
  \toprule
  \textbf{ID} & \textbf{Script} & \textbf{Optionnelle ?} \\
  \midrule
  01 & \texttt{01\_pre\_analyse.py} & Non \\
  02 & \texttt{02\_staging\_to\_bronze.py} & Non \\
  03 & \texttt{03\_bronze\_to\_silver.py} & Non \\
  04 & \texttt{04\_validation\_silver.py} & \textbf{Oui} \\
  05 & \texttt{05\_silver\_to\_gold.py} & Non \\
  06 & \texttt{06\_compiler\_panel.py} & Non \\
  06b & \texttt{06b\_imputer\_grades.py} & Non \\
  07 & \texttt{07\_calcul\_indicateur.py} & Non \\
  08 & \texttt{08\_creation\_fichier\_excel\_avec\_sommaire.py} & Non \\
  \bottomrule
\end{tabularx}

\section{Arguments en ligne de commande}

\begin{codebox}[]
python orchestrateur.py                          # pipeline complet
python orchestrateur.py --depuis 03             # reprend à partir de 03
python orchestrateur.py --etapes 06b 07 08      # étapes spécifiques
python orchestrateur.py --dry-run               # affiche le plan sans exécuter
python orchestrateur.py --depuis 07 --dry-run   # plan partiel
python orchestrateur.py --repertoire /mon/chemin # scripts dans un autre dossier
\end{codebox}

\begin{tabularx}{\textwidth}{lX}
  \toprule
  \textbf{Argument} & \textbf{Description} \\
  \midrule
  \texttt{--depuis ID} &
    Reprend à partir de l'étape ID (incluse). IDs valides :
    01, 02, 03, 04, 05, 06, 06b, 07, 08 \\
  \texttt{--etapes ID...} &
    Lance uniquement les étapes listées (ordre du pipeline préservé) \\
  \texttt{--dry-run} &
    Affiche le plan d'exécution sans lancer aucun script \\
  \texttt{--repertoire DIR} &
    Répertoire contenant les scripts.
    Défaut : dossier de l'orchestrateur (\texttt{\_\_file\_\_}) \\
  \bottomrule
\end{tabularx}

\section{Gestion des échecs}

Pour chaque étape, deux types d'issues distingués :

\begin{itemize}
  \item \textbf{ÉCHEC} : le script s'est terminé avec un code retour non nul
        (erreur Python non interceptée, sys.exit(1), etc.)
  \item \textbf{EXCEPTION} : le sous-processus n'a pas pu être lancé (fichier
        script introuvable, permission refusée, etc.)
\end{itemize}

Comportement selon le type d'étape :

\begin{tabularx}{\textwidth}{llX}
  \toprule
  \textbf{Type} & \textbf{Issue} & \textbf{Comportement} \\
  \midrule
  Obligatoire & ÉCHEC & Arrêt du pipeline \\
  Obligatoire & EXCEPTION & Arrêt du pipeline \\
  Optionnelle (04) & ÉCHEC & Pipeline continue \\
  Optionnelle (04) & EXCEPTION & Pipeline continue \\
  \bottomrule
\end{tabularx}

\section{Résumé final}

À la fin de l'exécution, l'orchestrateur affiche :
\begin{itemize}[nosep]
  \item Le statut (succès / échec) de chaque étape avec sa durée
  \item Les étapes non lancées (suite à un arrêt anticipé)
  \item La durée totale du pipeline
  \item La commande pour reprendre depuis le premier échec :
        \texttt{python orchestrateur.py --depuis \{id\_premier\_echec\}}
\end{itemize}

% ═════════════════════════════════════════════════════════════
% CHAPITRE 6 — GUIDE DE DÉMARRAGE RAPIDE
% ═════════════════════════════════════════════════════════════
\chapter{Guide de démarrage rapide}

\section{Prérequis}

\begin{itemize}
  \item Python 3.10+ avec les dépendances installées :
        \texttt{pip install polars fastexcel boto3 python-dotenv openpyxl}
  \item Serveur MinIO accessible (local ou JupyterHub)
  \item Fichiers sources disponibles localement :
    \begin{itemize}[nosep]
      \item Dossier des fichiers Excel mensuels
      \item \texttt{FICHIER\_ANSTAT\_CODE.xlsx}
      \item \texttt{CORRESPONDANCE\_GRILLE\_EMPLOI.xlsx}
    \end{itemize}
\end{itemize}

\section{Étapes d'exécution complète}

\textbf{Étape 1 --- Configurer les chemins locaux dans 00.}

Ouvrir \texttt{00\_upload\_staging.py} et modifier les trois variables :

\begin{codebox}[]
DOSSIER\_EXCEL          = "chemin/vers/fichiers\_excel/"
FICHIER\_ANSTAT         = "chemin/vers/FICHIER\_ANSTAT\_CODE.xlsx"
FICHIER\_CORRESPONDANCE = "chemin/vers/CORRESPONDANCE\_GRILLE\_EMPLOI.xlsx"
\end{codebox}

\textbf{Étape 2 --- Configurer la connexion MinIO.}

Dans chaque script, vérifier les valeurs de \texttt{MINIO\_ENDPOINT},
\texttt{MINIO\_ACCESS\_KEY}, \texttt{MINIO\_SECRET\_KEY}. Décommenter le
bloc JupyterHub si nécessaire.

\textbf{Étape 3 --- Déposer les fichiers sources (une seule fois).}

\begin{codebox}[]
python 00\_upload\_staging.py
\end{codebox}

Vérifier que les compteurs \texttt{OK / erreurs} sont cohérents.

\textbf{Étape 4 --- (Facultatif) Pré-analyser la structure.}

\begin{codebox}[]
python 01\_pre\_analyse.py
\end{codebox}

Consulter les rapports dans \bucket{staging/panel\_admin/pre\_analyse/}
pour identifier les variations de colonnes entre périodes.

\textbf{Étape 5 --- Lancer le pipeline complet.}

\begin{codebox}[]
python orchestrateur.py
\end{codebox}

Ou, si l'étape 01 a déjà été exécutée manuellement :

\begin{codebox}[]
python orchestrateur.py --depuis 02
\end{codebox}

\textbf{Étape 6 --- Vérifier la qualité (Silver).}

Consulter le rapport de matching dans
\bucket{staging/panel\_admin/validation/taux\_matching\_anstat\_par\_periode.csv}.
Un taux de matching faible sur \texttt{CODE\_ORGANISME} indique des libellés
non reconnus dans le fichier ANSTAT.

\textbf{Étape 7 --- Vérifier le panel Gold.}

Consulter les fichiers dans \bucket{gold/panel\_admin/} et les statistiques
affichées par \texttt{06\_compiler\_panel.py}.

\textbf{Étape 8 --- Configurer les paramètres d'indicateurs.}

Ouvrir \texttt{07\_calcul\_indicateur.py} et ajuster :
\begin{codebox}[]
MODE\_SALAIRE\_BRUT = True   # ou False pour revenu salarial
PERIODE\_RECENTE   = True   # ou False pour la série complète
INCLURE\_ZEROS     = True   # ou False
SEUIL\_WINS        = 0.01   # ou None pour désactiver
\end{codebox}

\textbf{Étape 9 --- Relancer les étapes 07 et 08 seulement.}

\begin{codebox}[]
python orchestrateur.py --etapes 07 08
\end{codebox}

\textbf{Étape 10 --- Récupérer le fichier Excel final.}

Télécharger depuis
\bucket{staging/panel\_admin/exports\_gold/indicateurs\_*\_SOMMAIRE.xlsx}.

\section{Erreurs fréquentes}

\begin{tabularx}{\textwidth}{lX}
  \toprule
  \textbf{Erreur} & \textbf{Cause et solution} \\
  \midrule
  \texttt{Fichier ANSTAT introuvable dans s3://staging/...} &
    Le fichier ANSTAT n'a pas été uploadé. Vérifier le chemin dans
    \texttt{FICHIER\_ANSTAT} et relancer \texttt{00\_upload\_staging.py}. \\
  \texttt{Colonnes GRADE et/ou GRADE\_SOURCE absentes du Gold} &
    L'étape 06b n'a pas encore été exécutée.
    Lancer \texttt{python orchestrateur.py --depuis 06b}. \\
  \texttt{Fichier introuvable : s3://staging/.../indicateurs\_*.xlsx} &
    Les paramètres \texttt{MODE\_SALAIRE\_BRUT} / \texttt{PERIODE\_RECENTE} /
    \texttt{INCLURE\_ZEROS} dans 08 diffèrent de ceux utilisés dans 07.
    Les aligner. \\
  \texttt{Aucune donnée pour les années [2024, 2025]} &
    Les données ne couvrent pas les années sélectionnées.
    Passer \texttt{PERIODE\_RECENTE = False}. \\
  \texttt{Taux matching barème $<$ 20\%} &
    Peu d'emplois reconnus dans \texttt{CORRESPONDANCE\_GRILLE\_EMPLOI.xlsx}.
    Vérifier les libellés d'emploi dans ce fichier. \\
  \texttt{ConnectionRefusedError / EndpointResolutionError} &
    MinIO inaccessible. Vérifier \texttt{MINIO\_ENDPOINT} (local vs JupyterHub). \\
  \texttt{NoSuchBucket} &
    Buckets non créés. Relancer \texttt{00\_upload\_staging.py} qui les crée
    automatiquement. \\
  \bottomrule
\end{tabularx}

\section{Relancer depuis un point d'échec}

Si l'orchestrateur s'arrête sur l'étape 03 :
\begin{codebox}[]
# Corriger le problème, puis :
python orchestrateur.py --depuis 03
\end{codebox}

Si seuls les indicateurs doivent être recalculés (paramètres modifiés) :
\begin{codebox}[]
python orchestrateur.py --etapes 07 08
\end{codebox}

% ═════════════════════════════════════════════════════════════
% CHAPITRE 7 — RÉFÉRENCE RAPIDE
% ═════════════════════════════════════════════════════════════
\chapter{Référence rapide}

\section{Séquence complète des scripts}

\begin{tabularx}{\textwidth}{llX}
  \toprule
  \textbf{Script} & \textbf{Buckets} & \textbf{Résultat} \\
  \midrule
  \texttt{00\_upload\_staging.py} &
    local $\rightarrow$ staging &
    Fichiers Excel + référentiels dans staging \\
  \texttt{01\_pre\_analyse.py} &
    staging $\rightarrow$ staging &
    Rapports CSV de structure des colonnes \\
  \texttt{02\_staging\_to\_bronze.py} &
    staging $\rightarrow$ bronze &
    \texttt{MMYYYY.parquet} (tout String) \\
  \texttt{03\_bronze\_to\_silver.py} &
    bronze + staging $\rightarrow$ silver &
    \texttt{YYYY.parquet} (enrichi, typé) \\
  \texttt{04\_validation\_silver.py} &
    silver $\rightarrow$ staging &
    Rapports qualité CSV (optionnel) \\
  \texttt{05\_silver\_to\_gold.py} &
    silver $\rightarrow$ gold + staging &
    \texttt{masse\_salariale.parquet},
    \texttt{effectifs.parquet}, CSV exports \\
  \texttt{06\_compiler\_panel.py} &
    silver $\rightarrow$ gold &
    \texttt{panel\_YYYY.parquet},
    \texttt{panel\_complet.parquet} \\
  \texttt{06b\_imputer\_grades.py} &
    gold $\rightarrow$ gold &
    Mise à jour GRADE + GRADE\_SOURCE dans tous les fichiers Gold \\
  \texttt{07\_calcul\_indicateur.py} &
    gold + staging $\rightarrow$ staging &
    \texttt{indicateurs\_*.xlsx} \\
  \texttt{08\_creation\_fichier\_excel\_avec\_sommaire.py} &
    staging $\rightarrow$ staging &
    \texttt{indicateurs\_*\_SOMMAIRE.xlsx} \\
  \bottomrule
\end{tabularx}

\section{Colonnes du Gold (panel)}

Le tableau suivant liste toutes les colonnes du panel Gold après exécution
complète du pipeline (étapes 06 + 06b) :

\begin{longtable}{lll}
  \toprule
  \textbf{Colonne} & \textbf{Type} & \textbf{Origine} \\
  \midrule
  \endhead
  \texttt{matricule} & Utf8 & F1 \\
  \texttt{nom} & Utf8 & F1 \\
  \texttt{sexe} & Utf8 & F1 \\
  \texttt{date\_naissance} & Utf8 & F1 \\
  \texttt{situation\_matrimoniale} & Utf8 & F1 \\
  \texttt{nombre\_enfant} & Utf8 & F1 \\
  \texttt{situation\_normalisee} & Utf8 & 03 (calcul) \\
  \texttt{CODE\_ORGANISME} & Utf8 & ANSTAT (03) \\
  \texttt{organisme} & Utf8 & F1 \\
  \texttt{CODE\_EMPLOI} & Utf8 & ANSTAT (03) \\
  \texttt{emploi} & Utf8 & F1 \\
  \texttt{Code\_CITP} & Utf8 & ANSTAT (03) \\
  \texttt{Metier\_CITP} & Utf8 & ANSTAT (03) \\
  \texttt{CODE\_GRADE} & Utf8 & ANSTAT (03) \\
  \texttt{grade} & Utf8 & F1 \\
  \texttt{statut\_fonctionnaire} & Utf8 & F1 \\
  \texttt{GRADE} & Utf8 & \textbf{06b} (OBSERVE/IMPUTE/NF) \\
  \texttt{GRADE\_SOURCE} & Utf8 & \textbf{06b} \\
  \texttt{CODE\_FONCTION} & Utf8 & ANSTAT (03) \\
  \texttt{fonction} & Utf8 & F1 \\
  \texttt{CODE\_SERVICE} & Utf8 & ANSTAT (03) \\
  \texttt{service} & Utf8 & F1 \\
  \texttt{CODE\_AFFECTATION} & Utf8 & ANSTAT (03) \\
  \texttt{lieu\_affectation} & Utf8 & F1 \\
  \texttt{CODE\_POSTE} & Utf8 & ANSTAT (03) \\
  \texttt{prise\_service} & Utf8 & F1 \\
  \texttt{date\_retraite} & Utf8 & F1 \\
  \texttt{age\_retraite} & Utf8 & F1 \\
  \texttt{mois\_annee} & Utf8 & Nom de fichier (03) \\
  \texttt{annee} & Int32 & 06 (calcul) \\
  \texttt{mois} & Int32 & 06 (calcul) \\
  \texttt{montant\_brut} & Float64 & F1 \\
  \texttt{montant\_net} & Float64 & F1 \\
  \texttt{retenue\_pension} & Float64 & F1 \\
  \texttt{impot} & Float64 & F1 \\
  \texttt{charge\_patronale} & Float64 & F1 \\
  \texttt{fichier\_source} & Utf8 & 02 (métadonnée) \\
  \texttt{124\_F2} \ldots \texttt{N\_F2} & Float64 & F2 (via 02) \\
  \texttt{CODE\_RÉGION} & Utf8 & ANSTAT (03) \\
  \texttt{CODE\_SITE} & Utf8 & ANSTAT (03) \\
  \texttt{CODE\_DRB} & Utf8 & ANSTAT (03) \\
  \texttt{CODE\_DÉPARTEMENT} & Utf8 & ANSTAT (03) \\
  \texttt{MONTANT TRANSPORT} & Utf8 & ANSTAT (03) \\
  \texttt{CODE\_ÉTABLISSEMENT} & Utf8 & ANSTAT (03) \\
  \texttt{CODE\_ACCT} & Utf8 & ANSTAT (03) \\
  \texttt{GRADE\_ADMINISTRATIF\_ASSOCIE} & Utf8 & ANSTAT (03) \\
  \bottomrule
\end{longtable}

\section{Feuilles Excel produites par 07}

Voir section 4.9 (tableau des feuilles produites).

\section{Exemple de chargement du panel}

\begin{codebox}[Charger le panel en Python]
import polars as pl, io, boto3
from botocore.client import Config

s3 = boto3.client("s3", endpoint\_url="http://192.168.1.230:30137",
                  aws\_access\_key\_id="datalab-team",
                  aws\_secret\_access\_key="minio-datalabteam123",
                  config=Config(signature\_version="s3v4"),
                  region\_name="us-east-1", verify=False)

# Charger une seule année
buf = io.BytesIO()
s3.download\_fileobj("gold", "panel\_admin/panel\_2023.parquet", buf)
buf.seek(0)
panel = pl.read\_parquet(buf)

# Charger toutes les années
frames = []
for annee in range(2015, 2026):
    buf = io.BytesIO()
    s3.download\_fileobj("gold", f"panel\_admin/panel\_{annee}.parquet", buf)
    buf.seek(0)
    frames.append(pl.read\_parquet(buf))
panel = pl.concat(frames, how="diagonal")

# Exemple : brut moyen par organisme et mois
panel.group\_by(["annee", "mois", "CODE\_ORGANISME"]).agg(
    pl.col("montant\_brut").mean().alias("brut\_moyen"),
    pl.col("matricule").n\_unique().alias("nb\_agents"),
).sort(["annee", "mois"])
\end{codebox}

% ═════════════════════════════════════════════════════════════
% CHAPITRE 8 — GLOSSAIRE
% ═════════════════════════════════════════════════════════════
\chapter{Glossaire}

\begin{description}[leftmargin=3cm, labelwidth=2.8cm, style=nextline]

  \item[ANSTAT] Système national de codification administrative. Le fichier
  \texttt{FICHIER\_ANSTAT\_CODE.xlsx} contient les dictionnaires de codes pour
  les organismes, grades, emplois, affectations, services, fonctions et postes.

  \item[Barème salarial] Grille de rémunération associée à un corps d'emploi
  (ex: Barème Général, ENS SUP, MAGISTRAT). La table
  \texttt{CORRESPONDANCE\_GRILLE\_EMPLOI.xlsx} fait correspondre chaque libellé
  d'emploi à son barème.

  \item[Bronze] Couche de données brutes : données issues des fichiers Excel,
  converties en Parquet sans aucune transformation métier, tout en
  \texttt{String}. Philosophie : ``Bronze = fidèle à la source''.

  \item[CITP] Classification Internationale Type des Professions.
  Le code CITP est rattaché via la table ANSTAT (\texttt{CODE\_EMPLOI}).
  Le premier chiffre identifie le grand groupe (1 = cadres, 2 = professions
  intellectuelles, etc.).

  \item[CLE\_UNIQUE] Identifiant composite formé par la concaténation du
  matricule et du code organisme (\texttt{matricule||code\_organisme}).
  Utilisé pour la jointure entre la Feuille 1 et la Feuille 2 des fichiers
  Excel.

  \item[CORRESPONDANCE\_GRILLE\_EMPLOI.xlsx] \textbf{(Nouveau)} Référentiel
  de correspondance entre libellé d'emploi et barème salarial, utilisé pour
  le fuzzy join dans l'étape 07.

  \item[F1 / F2] Les deux feuilles de chaque fichier Excel mensuel.
  F1 = données agents (identité, situation, montants).
  F2 = codes de rémunération élémentaires (composantes du salaire brut,
  identifiées par codes numériques).

  \item[Fuzzy join] Jointure sur libellé normalisé (après suppression des
  accents, ponctuation, passage en majuscules). Utilisée pour rattacher les
  codes ANSTAT et les barèmes malgré les variations orthographiques.

  \item[Gold] Couche de données finales prêtes à l'analyse : agrégats
  (masse salariale, effectifs) et panel individu$\times$période enrichi.

  \item[GRADE] Colonne ajoutée par l'étape 06b. Contient le grade fonctionnaire
  extrait ou imputé (ex: A3, B2, C1). Distinct de \texttt{grade} (libellé brut
  issu de F1).

  \item[GRADE\_SOURCE] Traçabilité du grade : \texttt{OBSERVE} (extrait
  directement), \texttt{IMPUTE} (inféré statistiquement par cascade emploi),
  \texttt{NF\_NON\_IMPUTABLE} (non-fonctionnaire ou emploi sans historique).

  \item[Imputation en cascade] Méthode d'imputation du grade en trois niveaux
  de priorité : P1 = mode du grade pour cet emploi ce mois précis,
  P2 = mode pour cet emploi cette année, P3 = mode pour cet emploi sur toute
  la série.

  \item[Medallion Architecture] Architecture en couches Bronze / Silver / Gold
  inspirée du \textit{Data Lakehouse}. Chaque couche améliore la qualité et la
  valeur des données.

  \item[MinIO] Serveur de stockage objet compatible S3, auto-hébergé.
  Accessible via le client boto3 de la même façon qu'AWS S3.

  \item[MMYYYY] Format des noms de fichiers mensuels (ex: \texttt{012015} =
  janvier 2015). Utilisé comme identifiant de période dans le pipeline.

  \item[MODELE\_ANSTAT] Feuille Excel au format hiérarchique présentant les
  statistiques de salaire brut (min/moyen/max) par secteur d'activité, dans
  la structure définie par l'ANSTAT pour la comparaison internationale.

  \item[Panel] Base de données longitudinale où chaque ligne représente un
  agent à un mois donné (format \textit{individu $\times$ période}).

  \item[Parquet] Format de fichier colonne binaire, compressé (Snappy dans ce
  pipeline), optimisé pour l'analyse analytique. Lecture rapide par Polars.

  \item[Polars] Bibliothèque Python de manipulation de DataFrames, alternative
  à pandas, avec API Lazy (scan/sink) permettant le streaming de grands
  volumes.

  \item[Silver] Couche de données enrichies : types corrects (Float64 pour les
  montants), situations normalisées, codes ANSTAT rattachés, lignes
  ``autre'' supprimées.

  \item[Staging] Couche d'atterrissage des fichiers sources. Zone de dépôt
  initiale et stockage des référentiels et rapports intermédiaires.

  \item[statut\_fonctionnaire] Colonne brute F1 contenant la description
  textuelle du grade (ex: ``Catégorie A3 -- Échelon 2''). La fonction
  \texttt{extraire\_grade()} en extrait le grade via regex.

  \item[Winsorisation] Écrêtage statistique des valeurs extrêmes. Dans ce
  pipeline, le REVENU\_TOTAL est écrêté à P1/P99 dans chaque cellule
  \texttt{mois $\times$ grade $\times$ CITP} (avec repli sur
  \texttt{mois $\times$ grade} si effectif insuffisant). Le Gold n'est pas
  modifié : la winsorisation est appliquée en mémoire uniquement.

\end{description}

\end{document}
